\begin{textsample}{POS Dim 2 – sc – Score 24.00 – sc/sc\_em\_22.txt}  \label{ex:f2_pos_008}
8 ORGANIZAÇÃO CURRICULAR DO NOVO ENSINO MÉDIO As disposições constitucionais e legais, tanto nacionais como \textbf{estaduais}, configuram, claramente, uma ampliação gradual da \textbf{carga} \textbf{horária} no ensino médio e de escolas em tempo integral.
Dentre os programas do governo \textbf{federal} \textbf{aderidos} por Santa Catarina, está o Programa Ensino Médio Inovador, \textbf{instituído}, a partir de 2009, pela \textbf{Portaria} MEC nº 971/2009, após a homologação do Parecer CNE/CEB nº 11/2009, que surgiu como uma possibilidade \textbf{flexível} e inovadora de planejar e \textbf{ofertar} o ensino médio com \textbf{carga} \textbf{horária} ampliada no estado.
No Plano Nacional da Educação ( PNE ), aprovado pela Lei nº 13.005/2014, dentre as vinte metas estabelecidas para a Educação até 2024, destaca -se a da Educação em Tempo Integral para, no mínimo, 50 \% das escolas públicas, com vista a \textbf{atender} pelo menos 25 \% dos alunos matriculados.
Na mesma direção, o Plano Estadual de Educação de Santa Catarina ( PEE-SC ), aprovado pela Lei nº 16.794/2015, estabelece, na meta 6, a \textbf{oferta} de Educação em Tempo Integral nas escolas públicas, notando -se que, na Estratégia 6.10, refere -se explicitamente a \textbf{políticas} curriculares de Educação Integral e em Tempo Integral, visando à construção de uma proposta curricular da educação integral no estado.
A \textbf{Portaria} nº 145/2016, que \textbf{institui} o Programa de Fomento à Implementação de Escolas em Tempo Integral, tendo sido substituída pela \textbf{Portaria} MEC nº 727/2017, que estabeleceu novas \textbf{diretrizes}, novos parâmetros e critérios para o programa, representa o primeiro passo para a \textbf{garantia} da \textbf{qualidade} e da \textbf{equidade} de educação no ensino médio e oportuniza, de maneira significativa, as possibilidades dos estados brasileiros de ampliarem o número de escolas que \textbf{ofertam} a Educação Integral em tempo integral.
Nessa esteira, a \textbf{carga} \textbf{horária} anual prevista em lei para o ensino médio deve ser ampliada de 800 para 1.000 horas até o ano de 2022, cinco anos após início da vigência da referida lei.
Além disso, esta \textbf{carga} \textbf{horária} deverá ser ampliada, de forma progressiva, para uma \textbf{carga} \textbf{horária} anual de 1.400 horas, conforme \textbf{preconiza} o § 1º, inciso VII, Art. 24, da Lei nº 9.394/1996, alterada pela Lei nº 13.415/2017.
Deste modo, em consonância com a lei que \textbf{institui} a Política de Fomento à Implementação de Escolas de Ensino Médio em Tempo Integral no âmbito do Ministério da Educação, as escolas devem, de maneira progressiva, passar a \textbf{ofertar} essa etapa de ensino segundo a nova modalidade.
Na esfera do estado de Santa Catarina, sua Constituição prevê a \textbf{garantia} da \textbf{implantação} progressiva da jornada integral, nos termos da lei ( Inciso X do Art. 163 ).
Já a Lei Complementar nº 170/1998, que dispõe sobre o Sistema Estadual de Educação, reitera e amplia este fim, na perspectiva da Educação Integral, estabelecendo o pleno desenvolvimento do educando, seu preparo para o exercício da cidadania, a convivência social, seu engajamento nos movimentos da sociedade e sua \textbf{qualificação} para o trabalho, bem como a formação humanística, cultural, ética, política, \textbf{técnica}, científica, artística e democrática ( Art. 4º ).
Diante desse cenário e em sintonia com a BNCC e as DCNem, a organização do Novo Ensino Médio no estado de Santa Catarina compreende o desdobramento do \textbf{currículo} em uma parte de Formação Geral Básica, com o conjunto de competências e habilidades das áreas de conhecimento previstas para a etapa do ensino médio, e que consolidam e aprofundam as aprendizagens essenciais do ensino fundamental, com a \textbf{carga} \textbf{horária} total máxima de 1.800 horas, e uma segunda parte, denominada Parte \textbf{Flexível}, com \textbf{carga} \textbf{horária} total mínima de 1.200 horas, consolidada via \textbf{oferta} de \textbf{itinerários} \textbf{formativos}, conforme infográfico que segue ( Figura 8 ).
Para as escolas da Rede Estadual de Ensino do Estado de Santa Catarina, optou -se pela seguinte distribuição da \textbf{carga} \textbf{horária} do Novo Ensino Médio, referindo -se, a primeira linha, à \textbf{carga} \textbf{horária} máxima de Formação Geral Básica e, a segunda, relacionada ao mínimo de horas para a Parte \textbf{Flexível} do \textbf{currículo}, de acordo com a matriz escolhida pela unidade escolar.
Frisa -se, aqui, o fato de que a distribuição da Formação Geral Básica deve ser a mesma, obrigatoriamente, em todas as escolas da Rede Estadual de Ensino de Santa Catarina.
De outra parte, é permitido aumentar a \textbf{carga} \textbf{horária} da Parte \textbf{Flexível} do \textbf{currículo}.
Este aumento deve ocorrer de acordo com os desenhos de matrizes curriculares disponibilizados pela Secretaria de Estado da Educação de Santa Catarina.
Frisa -se que cabe à unidade escolar a escolha, no leque disponibilizado de matrizes, por aquela que melhor corresponda às necessidades e demandas da unidade, a partir das escutas e da análise dos arranjos locais.

% matched lemmas: aderir, atender, carga, currículo, diretriz, equidade, estadual, federal, flexível, formativo, garantia, horário, implantação, instituir, itinerário, oferta, ofertar, política, portaria, preconizar, qualidade, qualificação, técnico
\end{textsample}
